% !TEX TS-program = pdflatex
% !TEX encoding = UTF-8 Unicode

% This file is a template using the "beamer" package to create slides for a talk or presentation
% - Giving a talk on some subject.
% - The talk is between 15min and 45min long.
% - Style is ornate.

% MODIFIED by Jonathan Kew, 2008-07-06
% The header comments and encoding in this file were modified for inclusion with TeXworks.
% The content is otherwise unchanged from the original distributed with the beamer package.

\documentclass{beamer}
\setbeamercovered{invisible}



% Copyright 2004 by Till Tantau <tantau@users.sourceforge.net>.
%
% In principle, this file can be redistributed and/or modified under
% the terms of the GNU Public License, version 2.
%
% However, this file is supposed to be a template to be modified
% for your own needs. For this reason, if you use this file as a
% template and not specifically distribute it as part of a another
% package/program, I grant the extra permission to freely copy and
% modify this file as you see fit and even to delete this copyright
% notice. 


\mode<presentation>
{
  \usetheme{Warsaw}
  % or ...

  % or whatever (possibly just delete it)
}


\usepackage[english]{babel}
% or whatever

\usepackage[utf8]{inputenc}
% or whatever

\usepackage{times}
\usepackage[T1]{fontenc}
% Or whatever. Note that the encoding and the font should match. If T1
% does not look nice, try deleting the line with the fontenc.

%%% MATH RELATED

\input{macros.tex}

\title{MATH 350-2 Advanced Calculus} 
\subtitle
{} % (optional)

\author[W.R. Casper] % (optional, use only with lots of authors)
{W.R. Casper}
% - Use the \inst{?} command only if the authors have different
%   affiliation.

\institute[California State University Fullerton] % (optional, but mostly needed)
{
  Department of Mathematics\\
  California State University Fullerton}
% - Use the \inst command only if there are several affiliations.
% - Keep it simple, no one is interested in your street address.

\subject{Talks}
% This is only inserted into the PDF information catalog. Can be left
% out. 



% If you have a file called "university-logo-filename.xxx", where xxx
% is a graphic format that can be processed by latex or pdflatex,
% resp., then you can add a logo as follows:

% \pgfdeclareimage[height=0.5cm]{university-logo}{university-logo-filename}
% \logo{\pgfuseimage{university-logo}}



% Delete this, if you do not want the table of contents to pop up at
% the beginning of each subsection:
\AtBeginSubsection[]
{
  \begin{frame}<beamer>{Outline}
    \tableofcontents[currentsection,currentsubsection]
  \end{frame}
}


% If you wish to uncover everything in a step-wise fashion, uncomment
% the following command: 

%\beamerdefaultoverlayspecification{<+->}


\begin{document}

\begin{frame}
  \titlepage
\end{frame}

\begin{frame}{Outline}
  \tableofcontents
  % You might wish to add the option [pausesections]
\end{frame}

% Since this a solution template for a generic talk, very little can
% be said about how it should be structured. However, the talk length
% of between 15min and 45min and the theme suggest that you stick to
% the following rules:  

% - Exactly two or three sections (other than the summary).
% - At *most* three subsections per section.
% - Talk about 30s to 2min per frame. So there should be between about
%   15 and 30 frames, all told.

\section{Real Analysis Lecture 12}


\subsection{Limits of sequences}

\begin{frame}{Definition of a limit}
\pause
Back in Calc. I:
\pause
$$\lim_{n\rightarrow\infty} x_n = L\quad\text{means}$$
\pause
that as $n$ gets arbitrarily large, the value of $x_n$ gets arbitrarily close to $L$.\\
\pause
\textbf{Formally:}
\pause
$$\forall\epsilon > 0,\quad\text{there exists $N\in \mathbb{Z}_+$}\ \ \text{such that}\ \ |x_n-L| < \epsilon\ \forall n\geq N.$$
\pause
\textbf{Metric space version:}
\pause
$$\forall\epsilon > 0,\quad\text{there exists $N\in \mathbb{Z}_+$}\ \ \text{such that}\ \ d(x_n,L) < \epsilon\ \forall n\geq N.$$
\pause
We say $\{x_n\}$ \textbf{converges} to $L$ and $L$ is the \textbf{limit of the sequence}.
\end{frame}

\begin{frame}{Example}
\pause
Consider the metric space $(M,d)$ defined by $M=\mathbb{R}$ with the Euclidean metric $d$.\\
\pause
The sequence $\{x_n\}$ defined by $x_n = 1/n$ converges to zero.\\
\pause
\begin{proof}
\pause
Let $\epsilon > 0$.\\
\pause
Choose $N\in\mathbb{Z}_+$ with $N > 1/\epsilon$.\\
\pause
Then $n>N$ implies
\pause
$$d(x_n,0) = |1/n-0| = 1/n < 1/N < \epsilon.$$
\pause
Since $\epsilon > 0$ was arbitrary, this proves $\{x_n\}$ converges to $0$.
\pause
\end{proof}
\end{frame}

\begin{frame}{Example}
Consider the metric space $(M,d)$ defined by $M=\mathbb{R}$ with the Euclidean metric $d$.\\
\pause
The sequence $\{x_n\}$ defined by $x_n = (n-1)/n$ converges to $1$.\\
\pause
\begin{proof}
\pause
Let $\epsilon > 0$.\\
\pause
Choose $N\in\mathbb{Z}_+$ with $N > 1/\epsilon$.\\
\pause
Then $n>N$ implies
$$d(x_n,1) = |(n-1)/n-1| = 1/n < 1/N < \epsilon.$$
\pause
Since $\epsilon > 0$ was arbitrary, this proves $\{x_n\}$ converges to $1$.
\end{proof}
\end{frame}

\begin{frame}{Example}
Consider the metric space $(M,d)$ defined by $M=\mathbb{R}$ with the {\color{red} discrete} metric $d$.\\
\pause
The sequence $\{x_n\}$ defined by $x_n = 1/n$ does not converge.\\
\pause
\begin{proof}
\pause
Suppose that it converges to a number $L$.\\
\pause
Then for all $\epsilon > 0$, there must exists an $N\in\mathbb{Z}_+$ such that $d(x_n,L) < \epsilon$ for all $n\geq N$.\\
\pause
Take $\epsilon = 1$ and choose $N\in\mathbb{Z}_+$ such that $d(x_n,L) < 1$ for all $n\geq N$.\\
\pause
With the discrete metric $d(x_n,L) < 1$ implies $x_n = L$.\\
\pause
This means $x_n=1/n = L$ for all $n\geq N$, which is a contradiction!
\end{proof}
\end{frame}

\begin{frame}{Uniqueness of limits}
\begin{thm}[Uniqueness of limits]
\pause
Let $(M,d)$ be a metric space.\\
\pause
If a sequence $\{x_n\}$ in $M$ converges to $L$ and converges to $T$, then $L=T$.
\end{thm}
\pause
\begin{proof}
\pause
Suppose that $\{x_n\}$ converges to $L$ and $T$.\\
\pause
Assume further that $L\neq T$ and let $\epsilon = d(L,T)/2$.\\
\pause
Then there exists $N_1$ such that $d(x_n,L) < \epsilon$ for $n \geq N_1$.\\
\pause
Likewise, there exists $N_2$ such that $d(x_n,T) < \epsilon$ for $n \geq N_2$.\\
\pause
Then for $n>\max\{N_1,N_2\}$,  the triangle inequality says
\pause
$$d(L,T) \leq d(L,x_n) + d(x_n,T) < 2\epsilon = d(L,T).$$
\pause
This is a contradiction.
\end{proof}
\end{frame}

\begin{frame}{Existence of limits}
\begin{thm}
If a sequence $\{x_n\}$ in a metric space $(M,d)$ converges to a value $L\in M$, then the range
$$X = \{x_1,x_2,\dots\}\subseteq M$$
is a bounded set and $L$ is an adherent point of $X$.
\end{thm}
\pause
When the range of a sequence is bounded, we call the sequence $\{x_n\}$ bounded.
\end{frame}

\begin{frame}{Existence of limits}
\begin{proof}
Suppose that $\lim_{n\rightarrow\infty} x_n = L$.\\
\pause
Let $\epsilon = 2024$.\\
\pause
Then there exists $N\in\mathbb{Z}_+$ with $d(x_n,L)\leq 2024$ for all $n\geq N$.\\
\pause
If we take
$$R = \max\{d(x_1,L),d(x_2,L),\dots, d(x_{N-1},L),2024\},$$
\pause
then $d(x_n,L)\leq R$ for all $N$.\\
\pause
Therefore $X\subseteq B_M(L,R)$ and in particular $X$ is bounded.
\end{proof}
\end{frame}

\begin{frame}{Existence of limits}
Sort of a converse of the previous theorem.
\pause
\begin{thm}
If $S$ is a subset of a metric space $(M,d)$ and $L\in M$ is an adherent point of $S$, then there exists a sequence $\{x_n\}$ of elements of $S$ which converges to $L$.
\end{thm}
\end{frame}

\begin{frame}{Existence of limits}
\begin{proof}
Suppose $L$ is an adherent point of $S$.\\
\pause
Then for all $n\in\mathbb{Z}_+$, the ball $B_M(x,\frac{1}{n})$ contains at least one element of $S$.\\
\pause
For each $n\in\mathbb{Z}_+$, choose $x_n\in S\cap B_M(x,\frac{1}{n})$.\\
\pause
We claim $\lim x_n = L$.\\
\pause
To see this, let $\epsilon > 0$.\\
\pause
Then choose $N\in\mathbb{Z}_+$ with $N>1/\epsilon$.\\
\pause
It follows that for all $n\geq N$
$$d(x_n,L) < 1/n < 1/N < \epsilon.$$
\pause
Since $\epsilon > 0$ was arbitrary, this proves convergence.
\end{proof}
\end{frame}


\end{document}


